\noindent
\begin{minipage}[t]{0.485\textwidth}

\noindent
\textbf{21.} \ $f(x) = \begin{cases} x + 3 & \text{nếu } x \le -1 \\ 2^x & \text{nếu } x > -1 \end{cases}$ \hfill \makebox[1.4cm][l]{$a = -1$}

\vspace{7pt}
\noindent
\textbf{22.} \ $f(x) = \begin{cases} \dfrac{x^2 - x}{x^2 - 1} & \text{nếu } x \neq 1 \\[6pt] 1 & \text{nếu } x = 1 \end{cases}$ \hfill \makebox[1.4cm][l]{$a = 1$}

\vspace{7pt}
\noindent
\textbf{23.} \ $f(x) = \begin{cases} \cos x & \text{nếu } x < 0 \\ 0 & \text{nếu } x = 0 \\ 1 - x^2 & \text{nếu } x > 0 \end{cases}$ \hfill \makebox[1.4cm][l]{$a = 0$}

\vspace{7pt}
\noindent
\textbf{24.} \ $f(x) = \begin{cases} \dfrac{2x^2 - 5x - 3}{x - 3} & \text{nếu } x \neq 3 \\[6pt] 6 & \text{nếu } x = 3 \end{cases}$ \hfill \makebox[1.4cm][l]{$a = 3$}

\vspace{6pt}
\noindent{\color{lightcyanline}\rule{\linewidth}{0.8pt}}
\vspace{5pt}

\noindent
\textbf{\color{stewartcyan}25--26}\enspace
\parbox[t]{0.82\linewidth}{%
  \textbf{(a)}~Chứng minh rằng $f$ có một điểm gián đoạn khử được tại $x = 3$.\\[2pt]
  \textbf{(b)}~Định nghĩa lại $f(3)$ để $f$ liên tục tại $x = 3$ (và do đó điểm gián đoạn được ``khử'').%
}

\vspace{6pt}
\noindent
\textbf{25.} \ $f(x) = \dfrac{x - 3}{x^2 - 9}$ \hfill \textbf{26.} \ $f(x) = \dfrac{x^2 - 7x + 12}{x - 3}$

\vspace{6pt}
\noindent{\color{lightcyanline}\rule{\linewidth}{0.8pt}}
\vspace{5pt}

\noindent
\textbf{\color{stewartcyan}27--34}\enspace Giải thích, bằng cách sử dụng các Định lý 4, 5, 7 và 9, tại sao hàm số sau đây liên tục tại mọi số thuộc tập xác định của nó. Nêu rõ tập xác định.

\vspace{6pt}
\noindent
\textbf{27.} \ $f(x) = \dfrac{x^2}{\sqrt{x^4 + 2}}$ \hfill \textbf{28.} \ $g(v) = \dfrac{3v - 1}{v^2 + 2v - 15}$

\vspace{7pt}
\noindent
\textbf{29.} \ $h(t) = \dfrac{\cos(t^2)}{1 - e^t}$

\vspace{6pt}
\noindent
\textbf{30.} \ $B(u) = \sqrt{3u - 2} + \sqrt[3]{2u - 3}$

\vspace{6pt}
\noindent
\textbf{31.} \ $L(v) = v \ln(1 - v^2)$ \hfill \textbf{32.} \ $f(t) = e^{-t^2}\ln(1 + t^2)$

\vspace{8pt}
\noindent
\textbf{33.} \ $M(x) = \sqrt{1 + \dfrac{1}{x}}$ \hfill \textbf{34.} \ $g(t) = \cos^{-1}(e^t - 1)$

\vspace{6pt}
\noindent{\color{lightcyanline}\rule{\linewidth}{0.8pt}}
\vspace{5pt}

\noindent
\textbf{\color{stewartcyan}35--38}\enspace Sử dụng tính liên tục để tính giới hạn.

\vspace{6pt}
\noindent
\textbf{35.} \ $\lim_{x \to 2} x\sqrt{20 - x^2}$ \hfill \textbf{36.} \ $\lim_{\theta \to \pi/2} \sin(\tan(\cos \theta))$

\vspace{9pt}
\noindent
\textbf{37.} \ $\lim_{x \to 1} \ln\left(\dfrac{5 - x^2}{1 + x}\right)$ \hfill \textbf{38.} \ $\lim_{x \to 4} 3^{\sqrt{x^2 - 2x - 4}}$

\vspace{6pt}
\noindent{\color{lightcyanline}\rule{\linewidth}{0.8pt}}
\vspace{5pt}

\noindent
\graphingcalcicon\ \textbf{\color{stewartcyan}39--40}\enspace Xác định các điểm gián đoạn của hàm số và minh họa bằng cách vẽ đồ thị.

\vspace{6pt}
\noindent
\textbf{39.} \ $f(x) = \dfrac{1}{\sqrt{1 - \sin x}}$ \hfill \textbf{40.} \ $y = \arctan(1/x)$

\end{minipage}%
\hfill
\begin{minipage}[t]{0.485\textwidth}

\noindent
\textbf{\color{stewartcyan}41--42}\enspace Chứng minh rằng $f$ liên tục trên $(-\infty, \infty)$.

\vspace{5pt}
\noindent
\textbf{41.} \ $f(x) = \begin{cases} 1 - x^2 & \text{nếu } x \le 1 \\ \ln x & \text{nếu } x > 1 \end{cases}$

\vspace{7pt}
\noindent
\textbf{42.} \ $f(x) = \begin{cases} \sin x & \text{nếu } x < \pi/4 \\ \cos x & \text{nếu } x \ge \pi/4 \end{cases}$

\vspace{9pt}
\noindent
\textbf{\color{stewartcyan}43--45}\enspace Tìm các số mà tại đó $f$ bị gián đoạn. Tại những số nào trong đó thì $f$ liên tục từ bên phải, từ bên trái, hay không bên nào? Vẽ phác đồ thị của $f$.

\vspace{5pt}
\noindent
\textbf{43.} \ $f(x) = \begin{cases} x^2 & \text{nếu } x < -1 \\ x & \text{nếu } -1 \le x < 1 \\ 1/x & \text{nếu } x \ge 1 \end{cases}$

\vspace{7pt}
\noindent
\textbf{44.} \ $f(x) = \begin{cases} 2^x & \text{nếu } x \le 1 \\ 3 - x & \text{nếu } 1 < x \le 4 \\ \sqrt{x} & \text{nếu } x > 4 \end{cases}$

\vspace{7pt}
\noindent
\textbf{45.} \ $f(x) = \begin{cases} x + 2 & \text{nếu } x < 0 \\ e^x & \text{nếu } 0 \le x \le 1 \\ 2 - x & \text{nếu } x > 1 \end{cases}$

\vspace{9pt}
\noindent
\textbf{46.} \ Lực hấp dẫn do Trái Đất tác dụng lên một vật có khối lượng đơn vị ở khoảng cách $r$ tính từ tâm Trái Đất là
\[
F(r) = \begin{cases}
\dfrac{GMr}{R^3} & \text{nếu } r < R \\[9pt]
\dfrac{GM}{r^2} & \text{nếu } r \ge R
\end{cases}
\]
trong đó $M$ là khối lượng của Trái Đất, $R$ là bán kính của nó, và $G$ là hằng số hấp dẫn. Hàm số $F$ có phải là một hàm liên tục của $r$ không?

\vspace{7pt}
\noindent
\textbf{47.} \ Với giá trị nào của hằng số $c$ thì hàm số $f$ liên tục trên $(-\infty, \infty)$?
\[
f(x) = \begin{cases}
cx^2 + 2x & \text{nếu } x < 2 \\[2pt]
x^3 - cx & \text{nếu } x \ge 2
\end{cases}
\]

\vspace{6pt}
\noindent
\textbf{48.} \ Tìm các giá trị của $a$ và $b$ để $f$ liên tục tại mọi điểm.
\[
f(x) = \begin{cases}
\dfrac{x^2 - 4}{x - 2} & \text{nếu } x < 2 \\[9pt]
ax^2 - bx + 3 & \text{nếu } 2 \le x < 3 \\[2pt]
2x - a + b & \text{nếu } x \ge 3
\end{cases}
\]

\vspace{6pt}
\noindent
\textbf{49.} \ Giả sử $f$ và $g$ là các hàm số liên tục sao cho $g(2) = 6$ và $\lim_{x \to 2} [3f(x) + f(x)g(x)] = 36$. Tìm $f(2)$.

\vspace{6pt}
\noindent
\textbf{50.} \ Cho $f(x) = 1/x$ và $g(x) = 1/x^2$.\\
\hspace*{1.2em}\textbf{(a)}~Tìm $(f \circ g)(x)$.\\
\hspace*{1.2em}\textbf{(b)}~Hàm số $f \circ g$ có liên tục tại mọi điểm không? Giải thích.

\end{minipage}

\vfill
\begin{center}
  \tiny\color{gray}
  Bản quyền thuộc Cengage Learning 2021. Mọi quyền được bảo lưu. Không được sao chép, quét hoặc nhân bản một phần hay toàn bộ. WCN 02-200-203\\
  Đánh giá của ban biên tập cho thấy rằng mọi nội dung bị lược bỏ không ảnh hưởng đáng kể đến trải nghiệm học tập tổng thể. Cengage Learning bảo lưu quyền xóa nội dung bổ sung vào bất kỳ lúc nào nếu các hạn chế về quyền sau đó yêu cầu.
\end{center}
